\section{What We Found}\label{sec:findings}

Table~\ref{tab:findings} summarizes what the formalization found. One
item is a gap in a proof. The others are misprints and small slips:
statements that are false as printed but obvious to correct, wrong
citations inside derivations, and a few conventions the papers use but
do not state. Appendix~\ref{sec:appendix-errata} describes those one by
one, and says how each was found. Here we describe the gap, and then
the errors we made ourselves.

\begin{table}
\caption{Errors found in the sources. ``Statement'' means a printed
statement that is false; ``derivation'' a wrong step or citation inside
a proof whose conclusion is right.}
\label{tab:findings}
\small
\begin{tabularx}{\linewidth}{@{}l Y l@{}}
\toprule
Source & Finding & Kind \\
\midrule
\citet{fang2026hadamard} & Lemmas A.16 and A.17 are false; Main Case 9 of the Main Lemma fails on a state in dimension 9 & gap \\
 & Two wrong citations in Lemma A.19(d); one needed relation is stated nowhere; Prop.~A.29 repeats Prop.~A.22 & derivation \\
\citet{makary2021real} & 11 of the 137 typed relations of Appendix~B are false as printed & statement \\
 & Figure~10 lacks an ``=''; Definition~5.3 omits a case that rules 5.1--5.4 use & statement \\
 & Supplement: 6 derivations reach a false conclusion, 5 misprinted rows, 4 wrong axiom labels & derivation \\
\citet{selinger2015clifford} & Supplement: two steps labelled (11) are (19); Lemma 1.3 cites (C10) for (C12) & derivation \\
\citet{li2021on} & Relation (7d) of Table~2 has its two right-hand sides exchanged & statement \\
\citet{blake2026simpler} & The scalar of the qutrit rule (E) in Figure~4 is $\omega_{12}^7$; it must be $\omega_{12}$ & statement \\
 & One separating interpretation is not a model of a three-wire rule & derivation \\
\citet{clement2026realcliffordch} & Its Theorem 4.4 is Theorem 4.14 of \citet{fang2026hadamard}, whose proof has the gap above & dependency \\
 & (36)/(37) hold with $\mathcal{E}^{x}$ where $\mathcal{E}^{1-x}$ is printed; distinctness conventions of (30), (34), (42) unstated & statement \\
\bottomrule
\end{tabularx}
\end{table}

\subsection{A gap in the completeness proof of Hadamard-\texorpdfstring{$\Pi$}{Pi}}\label{sec:gap}

Hadamard-$\Pi$~\cite{fang2026hadamard} proves the completeness of its
Figure~6 (Theorem~4.14) by the method of
Section~\ref{sec:by-levels}. The level of a matrix is $(j,k,\ell)$: the
pivot column $j$, its least denominator exponent $k$, and the number
$\ell$ of entries of $\sqrt2^k$ times the column whose integer part is
odd. The odd entries come in two residue classes modulo 2, $1$ and
$1+\sqrt2$, each of even size. Algorithm~1 applies $H$ to the first two
odd entries of one class, which makes both even. The Main Lemma
(Lemma~A.7) is split into Main Cases by the generator $G$ and the
syllable $N$. Main Case~9, where $G=H_{[1,2]}$ and $N=H_{[1,i_2]}$, takes
nine pages in the extended version. In one of its subcases it applies
Lemma~A.19, which gives a path between two states $t_2$ and $t_3$ that
stays below a promised level. In the case $\ell=4$, sub-case (2) of
Lemma~A.19, items (c) and (d) rely on Lemmas~A.16 and A.17.

These two lemmas are false. Lemma~A.16 assumes that $M$ is at level
$(j,k,4)$ with $k\ge2$ and that the first four entries $v$ of column $j$
satisfy $\sqrt2^k v \equiv (1+\sqrt2,1,1+\sqrt2,1)$ or
$(1,1+\sqrt2,1,1+\sqrt2) \pmod 2$. It states that if
$\sqrt2^{k-1}H_{[1,4]}H_{[2,3]}H_{[2,4]}H_{[2,3]}H_{[1,4]}v$ has
entries in $\Z[\sqrt2]$, then so does
$\sqrt2^{k-1}H_{[1,3]}H_{[2,4]}H_{[1,2]}H_{[1,3]}H_{[2,4]}v$. Take
$n=j=6$, $k=4$ and the unit vector with
\[
\sqrt2^4\cdot M_{\cdot,6} = (-1-\sqrt2,\; 1,\; -1-\sqrt2,\; 1,\; \sqrt2,\; -2+\sqrt2).
\]
Some matrix of $O_6(\Z[\frac1{\sqrt2}])$ has this last column, because
the generators act transitively on unit columns. Its level is
$(6,4,4)$. The first word gives $\frac1{\sqrt2^{3}}(0,0,-1-\sqrt2,-1)$,
which has the required form. The second gives
$\frac1{\sqrt2^{3}}(-\frac{\sqrt2}2,-1-\frac{\sqrt2}2,-\frac{\sqrt2}2,-1-\frac{\sqrt2}2)$,
which does not. The proof of the lemma writes the entries as
$a_i+\sqrt2 b_i$ and compares the $\sqrt2$-parts of the norm equation
(page~38 of the extended version). There the constant terms of
$(2a_1+1)(2b_1+1)$ and $(2a_3+1)(2b_3+1)$ contribute $4$, so the
bracket that the proof reduces modulo~2 is missing a $+1$. With the
$+1$, the case $a_3+a_4$ odd is no longer contradictory, and the
vector above lies in it. Lemma~A.17 has the same hypotheses with
$H_{[3,4]}$ in place of $H_{[2,4]}$, ends its proof by appeal to the
argument of Lemma~A.16, and fails in the same way, in the case
$a_3+a_4$ even. A counterexample is the column
$(1+\sqrt2,-1,1+\sqrt2,1,-\sqrt2,2-\sqrt2)/\sqrt2^4$.

A false lemma does not make the Main Lemma false, since Lemma~A.19 only
needs some path, and there might be one. So we looked for a state on
which the proof itself fails. We sampled random states and followed the
case analysis of the Main Lemma on each of them. Such states are rare:
for $n=7$ and $n=8$ there were none among about 3{,}000 and 9{,}000
instances of the relevant sub-case, and for $n=10$ one among about 400
states that reach the step. One of them is the state $s$ with
\[
\sqrt2^7\cdot s_{\cdot,9} = (1+3\sqrt2,\; 1,\; -3\sqrt2,\; 1+3\sqrt2,\; 2,\; 5,\; -2+\sqrt2,\; -4+2\sqrt2,\; -2-2\sqrt2)
\]
at level $(9,7,4)$. Algorithm~1 takes $N=H_{[1,4]}$, and the proof of
Main Case~9 goes to case (2)(b)(i), then to its second bullet, then to
Lemma~A.19 with $t_2=H_{[2,6]}H_{[1,4]}s$ at level $(9,6,4)$ and
$(a,b,c,d)=(1,2,4,6)$. Rows 1 and 2 of $t_2$ have different residues,
which is sub-case (2), and only item (d) of Lemma~A.15 holds. Main
Case~9 needs a path from $t_2$ to $t_3$ strictly below $(9,7,4)$.
Every one of the eight diagrams of Lemma~A.19, started at $t_2$ with
these labels, reaches $(9,7,4)$ or $(9,8,4)$. So the proof of the Main
Lemma does not cover this state, and the proof of Theorem~4.14 is
incomplete.

In terms of states, what is left open is the edge given by $H$ on two
odd entries of different classes, when exactly four entries of the
pivot column are odd. Cl\'ement's 2019 thesis~\cite{clement2019thesis}
presents the same group with the same generators by its own
relations (1)--(21), with a completeness proof by the same method, and
his Case~3.4 is exactly this edge. We did two things with it.

First, we proved that his relations and those of Figure~6 generate the
same congruence in every dimension. His (1)--(18) follow from
(a1)--(d2), his (19) is a relation the paper derives from (d3), and
conversely. His (20) and (21) are relations on four and six indices
that are equivalent to (d3) and (d4) given the others; we derived each
direction by a generated chain of rewrites on the first indices, and
carried the result to any indices by an increasing embedding. So a
completeness proof for either set is one for both.

Second, we formalized his Case~3.4 and used it for the edge that
Hadamard-$\Pi$ leaves open. On the state $s$ above, his diagram~(32), on
rows 1, 2, 4, 6 and the further rows 3 and 5, gives a path that never
rises above the level of $s$. The general argument is a case analysis
over the residues of four entries, extended in some sub-cases by up to
three further rows of given depth and class, and closed by his
diagrams~(20) and~(30). We wrote it as a decision tree that Agda checks
by evaluation (Section~\ref{sec:trees}). With it, the relations of
Figure~6 present $O_n(\Z[\frac1{\sqrt2}])$ in every dimension.

After that we formalized the rest of the thesis: his version of the
algorithm, in which every syllable is a single generator, his
Lemma~4.4 for the basic generators $H_{[0,1]}$, $(-1)_{[j]}$ and
$X_{[j,j+1]}$ in all its cases and subcases, and his Theorem~4.1. Two
steps the thesis leaves implicit had to be supplied: the reduction from
all generators to the basic ones, by conjugation with exchanges that
keep a state at or below its level, and the induction itself. So the
theorem now has two machine-checked proofs, one along Hadamard-$\Pi$
with Cl\'ement's case, and one along Cl\'ement.

Cl\'ement's 2026 paper on real Clifford+$CH$
circuits~\cite{clement2026realcliffordch} imports Theorem~4.14 of
Hadamard-$\Pi$ as its own Theorem~4.4, and uses it twice: for circuits
on two qubits, and, through its Theorem~4.10, for three qubits and
more. Its main theorem therefore rested on a theorem whose published
proof is incomplete. Section~\ref{sec:encoding} explains how we connect
the two developments so that its completeness is proved with no
imported statement.

The two counterexamples and the state $s$ are given above, and all
three are small enough to check by hand; the errata will be provided
as supplementary material.

\subsection{Our own errors}\label{sec:own-errors}

The formalization found errors in our transcriptions as well, and some
of them are instructive.

\paragraph{A false hypothesis.} The real Clifford+$CH$ development was
first written with Theorem~4.4 as a module parameter, stated as
completeness of the relations of Figure~7 for every word over the
generators $(-1)_{[a]}$, $X_{[a,b]}$, $H_{[a,b]}$ with indices in
$\mathsf{Fin}\,N$. That statement is false. The type of generators
admits $X_{[a,a]}$, which denotes the identity, but every rule has
distinct indices, so the number of such letters is invariant under the
congruence and $X_{[a,a]}$ is not related to the empty word. Every
theorem proved from the hypothesis was vacuous. We noticed only when we
tried to discharge it. The fix restricts completeness to proper words,
whose letters have distinct indices, as the paper intends; the encoding
was then shown to produce only proper words. A hypothesis that nobody
discharges is not checked by anything.

\paragraph{A bug invisible at small sizes.} Our transcription of the
swap networks of Cl\'ement's Definition~8.2 took two products over the
wires between $i$ and $j$ in the wrong order. On three qubits there is
at most one wire in between, so the two orders agree, and the decision
procedure that checks Figure~8 on three qubits could not see the
difference. On four qubits the transcribed (36) and (37) are unsound:
a simulation of the network's signed permutation in Python showed that
8 of the 48 relevant index tuples land on the wrong basis vectors, and
48 of 160 on five qubits. We found it when proving (65) at every width,
which needs exactly this permutation. This is our error, distinct from
the paper's misprint of $\mathcal{E}^{1-x}$ in the same equations
(Appendix~\ref{sec:misprints}).

\paragraph{A misread diagram.} A control dot of equation (261),
misread at screen resolution, cost a proof of a true but harder
statement before the one the paper means; we now read diagrams at
400~dpi.
